\section{An arithmetic envelope for exact memory at every phase}\label{sec:r71-arithmetic}
The spectral criterion of Theorem~\ref{thm:r69-spectral} determines exact stochastic memory by membership in the classical stochastic eigenvalue regions. A worst-phase bound and a matching estimate at a badly approximable fixed phase do not describe all fixed phases. In particular, excellent rational approximations can produce both unusually small memory and long intervening ranges of larger memory. We determine a single arithmetic envelope that controls the entire memory function, and derive its exact lower and upper logarithmic growth exponents for every irrational phase. The arithmetic is classical continued-fraction theory; the conclusions here concern its conversion into the least exact stochastic observer memory. We retain the distinction from the noise-dependent classical process bounds in \cite{Fanizza24} and from the classical irrationality-exponent identity in \cite{HoneVarona19}.

Throughout this section, $0<\epsilon<1$, $\alpha=1-\epsilon$, and
\[
 b_n=\tfrac12[1+\gamma\alpha^n\cos(2\pi nx+\varphi)],
 \qquad 0<\gamma\le1,\quad \sin(2\pi x)\ne0.
\]
The phase and nonzero contrast describe one prepared target, not information revealed to its observer. Write $m_x(\epsilon)=m(1-\epsilon,2\pi x)$ for its least live-label count, with a prescribed initial label, stochastic transitions and a decoder in $[0,1]$. Theorem~\ref{thm:r69-spectral} makes this independent of $\gamma,\varphi$. The actual squared-loss embedding of Theorem~\ref{thm:r67-embedding} forces exactly this probability response at the least physical cost $1/p$. Thus the conclusions below are about exact optimal implementation, not approximation of the target sequence or a new definition of loss.

For a positive integer $q$, put
\[
 e_x(q)=\frac{\|qx\|_{\mathbb R/\mathbb Z}}q
       =\min_{a\in\mathbb Z}|x-a/q|,
 \qquad
 \mathfrak A_x(\epsilon)
   =\min_{q\ge1}\max\{3,q,e_x(q)/\epsilon\}.
\]
The minimum is attained: its summand is at least $q$, while $q=1$ gives a finite upper. No global Diophantine lower hypothesis is imposed.

\begin{theorem}[Universal arithmetic envelope]\label{thm:r71-envelope}
For every nonreal phase and every $0<\epsilon<1$,
\begin{equation}\label{eq:r71-envelope}
 \tfrac12\mathfrak A_x(\epsilon)
 \le m_x(\epsilon)
 \le C_0\mathfrak A_x(\epsilon),\qquad C_0=1+4\pi^2.
\end{equation}
The constants are independent of the phase, contrast, preparation phase and damping. Each upper bound comes from an explicit stochastic circulant realization with a prescribed initial label and bounded decoder. The lower holds over all stochastic matrices and decoders, not only those circulants. The estimate is a constant-factor law for the full arithmetic envelope; it is not an assertion of an exact leading constant or a new classification of the stochastic spectral regions.
\end{theorem}
\begin{proof}
We use the polygon construction and the power-sector inequality in Lemma~\ref{lem:r69-arithmetic-tools}, whose positive-matrix proofs are included in Section~\ref{sec:r69-spectral}. Let $M=m_x(\epsilon)\ge3$. Among the $2M+1$ fractional parts of $0,x,\ldots,2Mx$, two lie in the same one of $2M$ intervals of length $1/(2M)$. Their difference supplies $1\le k\le2M$ with $\|kx\|\le1/(2M)$. If this norm is zero, using $q=k$ gives $\mathfrak A_x(\epsilon)\le2M$ directly.

Otherwise put $\delta=2\pi\|kx\|$ and $h=\pi/M$. Then $0<\delta\le h\le\pi/3$. For
$F=\cos\delta+\sin\delta\tan h$, one has $F\le\sec h\le2$ and
\[
 F-1\ge\tfrac34h\delta-\tfrac12\delta^2\ge h\delta/4.
\]
Here $\sin\delta\ge3\delta/4$, $\tan h\ge h$ and $1-\cos\delta\le\delta^2/2$ on the stated interval. The eigenvalue of the stochastic power $Q^k$ gives $(1-\epsilon)^kF\le1$. Bernoulli's inequality therefore yields
\[
 k\epsilon\ge1-(1-\epsilon)^k\ge1-1/F\ge h\delta/8.
\]
Consequently $e_x(k)/\epsilon=\delta/(2\pi k\epsilon)\le4M/\pi^2$. Since also $k\le2M$, the envelope is at most $2M$. This proves the lower.

For the upper fix any denominator $q$, choose its closest integer numerator $a$, and choose the smallest positive multiple $l$ of $q$ satisfying
\[
 l\ge3,\qquad l\ge4\pi^2e_x(q)/\epsilon.
\]
Then
$l<q+\max\{3,4\pi^2e_x(q)/\epsilon\}\le C_0\max\{3,q,e_x(q)/\epsilon\}$.
The rational angle $2\pi a/q$ is an $l$th-root angle. Thus the distance $\delta_l$ from $2\pi x$ to the nearest such root is at most $2\pi e_x(q)$. The directional polygon bound gives
\[
 1-R_l(2\pi x)\le2(\pi/l)\delta_l
             \le4\pi^2e_x(q)/l\le\epsilon.
\]
Hence $\alpha e^{2\pi ix}$ lies in the regular $l$-gon. The positive combination of its two adjacent roots and zero gives the stochastic matrix described in Lemma~\ref{lem:r69-arithmetic-tools}; its normalized eigenvector gives the bounded decoder from Theorem~\ref{thm:r69-spectral}. The original initial label suffices. Minimize over $q$ to conclude. All constructions act on the observer's private memory and do not add a physical observation or a call.
\end{proof}

For irrational $x$, $\mathfrak A_x(\epsilon)\to\infty$ as $\epsilon\downarrow0$: a bounded envelope would restrict $q$ to finitely many integers whose errors have a positive minimum. For a rational nonreal phase of reduced denominator $q$, the exact memory is eventually $q$ as $\epsilon\downarrow0$. Indeed $\alpha C_q+(1-\alpha)J$ gives the upper $q$, while a bounded sequence of smaller stochastic realizations would converge, by compactness, to one with the forbidden unit eigenvalue of order $q$. This rational boundary is different from assigning an irrationality exponent to a rational number. Zero contrast and real phases remain governed by their separately stated degeneracies.

\subsection{The arithmetic input, with proof}
For irrational $x$, changing it by an integer does not change any quantity above. Write its fractional part as $[0;a_1,a_2,\ldots]$, with convergents $p_n/q_n$. We use the standard initial values $p_{-1}=1$, $p_0=0$, $q_{-1}=0$, $q_0=1$ and the recursion
$p_{n+1}=a_{n+1}p_n+p_{n-1}$, $q_{n+1}=a_{n+1}q_n+q_{n-1}$. The possible initial repeated denominator is irrelevant below; all asymptotic statements use the eventually strictly increasing sequence.

\begin{lemma}[Convergent errors and the irrationality exponent]\label{lem:r71-cf}
Set $E_n=q_nx-p_n$, with $x$ reduced modulo one. Then
\begin{equation}\label{eq:r71-cf-error}
 \frac1{q_{n+1}+q_n}<|E_n|<\frac1{q_{n+1}},
 \qquad
 |qx-a|\ge|E_n|\quad(0<q<q_{n+1},\ a\in\mathbb Z).
\end{equation}
For all sufficiently large $n$, $e_x(q_n)=|E_n|/q_n$. Define
\[
 \mu(x)=\sup\{t:\ e_x(q)<q^{-t}\text{ for infinitely many }q\}.
\]
Then
\begin{equation}\label{eq:r71-mu}
 \mu(x)=1+\limsup_{n\to\infty}\frac{\log q_{n+1}}{\log q_n}
       \in[2,\infty].
\end{equation}
Moreover $x$ is badly approximable, in the sense that $e_x(q)\ge c/q^2$ for some $c>0$ and all $q$, if and only if the ratios $q_{n+1}/q_n$ are bounded.
\end{lemma}
\begin{proof}
The recursions give $p_nq_{n-1}-p_{n-1}q_n=(-1)^{n-1}$ by induction. If $t_{n+1}=[a_{n+1};a_{n+2},\ldots]$, repeated substitution in the continued fraction gives
\[
 x=\frac{p_nt_{n+1}+p_{n-1}}{q_nt_{n+1}+q_{n-1}},
 \quad |E_n|=\frac1{q_nt_{n+1}+q_{n-1}},
 \quad E_{n-1}=-t_{n+1}E_n.
\]
Since $a_{n+1}<t_{n+1}<a_{n+1}+1$, the denominator lies strictly between $q_{n+1}$ and $q_{n+1}+q_n$, proving the first pair of bounds. These errors eventually have modulus below $1/2$, so the numerator is the nearest integer.

For the second assertion, suppose $|qx-a|<|E_n|$. The determinant identity makes the two consecutive convergent vectors an integral basis. Write
$(a,q)=A(p_n,q_n)+B(p_{n-1},q_{n-1})$ with integers $A,B$.
Then $|A-Bt_{n+1}|<1$. If $B=0$, this forces $A=0$ and $q=0$. If $B\ge1$, it forces $A\ge Ba_{n+1}$ and hence $q\ge Bq_{n+1}\ge q_{n+1}$. If $B=-b\le-1$, it forces $A\le-ba_{n+1}$ and hence $q\le-bq_{n+1}<0$. These exhaust the possibilities and prove the asserted best-error inequality. No best-approximation assertion for arbitrary fitted rational data has been used.

Put $\Lambda=\limsup\log q_{n+1}/\log q_n\ge1$. The convergent upper $e_x(q_n)<1/(q_nq_{n+1})$ shows $\mu\ge1+\Lambda$, including the infinite case. If $s>\Lambda$ is finite, eventually $q_{n+1}\le q_n^s$. For $q_n\le q<q_{n+1}$, the second-kind inequality gives
\[
 e_x(q)>\frac1{q(q_{n+1}+q_n)}
        \ge\frac1{2q^{s+1}}.
\]
For every positive $\eta$, this excludes infinitely many errors below $q^{-(s+1+\eta)}$. Let $s\downarrow\Lambda$ and $\eta\downarrow0$ to prove \eqref{eq:r71-mu}.

If $q_{n+1}/q_n\le B$ eventually, the same interval argument gives $e_x(q)>1/[(B+1)q^2]$ there; reducing the positive constant handles the finitely many earlier denominators. Conversely $e_x(q_n)<1/(q_nq_{n+1})$ is incompatible with one positive lower constant when the ratios are unbounded. This proves the last equivalence. The identity \eqref{eq:r71-mu} is classical; it is also stated in the cited arithmetic literature \cite{HoneVarona19}, and is proved here to make the memory argument independent of that reference.
\end{proof}

\subsection{The complete logarithmic growth spectrum}
The envelope law retains more information than a single exponent. Its two logarithmic exponents nevertheless have a closed form for every irrational phase, including phases not known to satisfy a uniform approximation bound.

\begin{theorem}[Exact lower and upper memory exponents]\label{thm:r71-exponents}
For every irrational $x$, the exact stochastic memory obeys
\begin{align}
 \liminf_{\epsilon\downarrow0}
  \frac{\log m_x(\epsilon)}{\log(1/\epsilon)}
     &=\frac1{\mu(x)+1},\label{eq:r71-liminf}\\
 \limsup_{\epsilon\downarrow0}
  \frac{\log m_x(\epsilon)}{\log(1/\epsilon)}
     &=\frac{\mu(x)-1}{2\mu(x)-1}.\label{eq:r71-limsup}
\end{align}
For $\mu(x)=\infty$, the right sides mean $0$ and $1/2$. In particular the logarithmic exponent exists exactly when $\mu(x)=2$, and then it equals $1/3$. This does not imply a constant-factor cubic-root law for every such phase. The assertions describe one fixed phase while the noise tends to zero, not a worst-phase choice depending on the noise.
\end{theorem}
\begin{proof}
By Theorem~\ref{thm:r71-envelope}, multiplicative constants do not affect either logarithmic limit, so work with $\mathfrak A_x$.

First suppose $\mu<\infty$ and choose $t>\mu$. For all sufficiently large $q$, $e_x(q)\ge q^{-t}$. Therefore
\[
 \max\{q,e_x(q)/\epsilon\}
   \ge\max\{q,q^{-t}/\epsilon\}
   \ge\epsilon^{-1/(t+1)}.
\]
Every one of the finitely many remaining denominators has positive error, and its error term grows as $\epsilon^{-1}$, so those denominators do not invalidate this lower for sufficiently small $\epsilon$. This gives a lower limit at least $1/(\mu+1)$. Conversely, for $1<t<\mu$, infinitely many $q$ satisfy $e_x(q)<q^{-t}$. At $\epsilon=q^{-(t+1)}$, the envelope is at most $q$ once $q\ge3$. Taking $t\uparrow\mu$ proves \eqref{eq:r71-liminf}; for $\mu=\infty$ the same argument with arbitrarily large $t$, together with nonnegativity, gives zero.

For the upper limit with finite $\mu$, again choose $t>\mu$. Pigeonhole approximation at denominator bound $Q$ gives $1\le k\le Q$ with $e_x(k)\le1/(kQ)$. Such $k$ tend to infinity along $Q\to\infty$, because the phase is irrational. The eventual lower $e_x(k)\ge k^{-t}$ therefore forces $k\ge Q^{1/(t-1)}$. Hence
\[
 \mathfrak A_x(\epsilon)
 \le\max\{3,Q,\epsilon^{-1}Q^{-t/(t-1)}\}.
\]
Choose $Q=\lceil\epsilon^{-(t-1)/(2t-1)}\rceil$ and then let $t\downarrow\mu$. This bounds the upper limit by the right side of \eqref{eq:r71-limsup}. If $\mu=\infty$, use only $k\ge1$ and $Q=\lceil\epsilon^{-1/2}\rceil$ in the same pigeonhole argument to obtain the universal upper limit $1/2$.

To prove the reverse upper-limit bound, take consecutive large convergent denominators $Q=q_n$, $Q'=q_{n+1}$ and set
\begin{equation}\label{eq:r71-crossover}
 \epsilon_n=\frac1{Q(Q')^2}.
\end{equation}
For every rational $a/q$ with $0<q<Q'/2$, either $a/q=p_n/Q$, in which case its error is greater than $1/[Q(Q'+Q)]\ge1/(2QQ')$, or the determinant of the two integer pairs is nonzero and
\[
 |x-a/q|\ge\frac1{qQ}-|x-p_n/Q|
          >\frac1{qQ}-\frac1{QQ'}>\frac1{QQ'}.
\]
Thus all these denominators have $e_x(q)/\epsilon_n\ge Q'/2$. All other denominators have $q\ge Q'/2$. It follows that $\mathfrak A_x(\epsilon_n)\ge Q'/2$. On the other hand, denominator $Q'$ has
$e_x(Q')<1/(Q'q_{n+2})$, so
$e_x(Q')/\epsilon_n<Q Q'/q_{n+2}\le Q\le Q'$.
Consequently
\begin{equation}\label{eq:r71-crossover-envelope}
 Q'/2\le\mathfrak A_x(\epsilon_n)\le Q'.
\end{equation}
The corresponding logarithmic ratio is, up to a term tending to zero,
\[
 \frac{\log Q'}{\log Q+2\log Q'}
   =\frac{\log Q'/\log Q}{1+2\log Q'/\log Q}.
\]
The function on the right is increasing, continuous and tends to $1/2$ at infinity. Lemma~\ref{lem:r71-cf} proves precisely the remaining lower bound in \eqref{eq:r71-limsup}. Finally the two displayed exponents are equal exactly when $\mu=2$: solving the equality gives $\mu(\mu-2)=0$, and $\mu\ge2$.
\end{proof}

\begin{theorem}[Rigidity of the uniform cubic-root law]\label{thm:r71-cubic-rigidity}
For an irrational fixed phase, the uniform estimate
\[
 m_x(\epsilon)=\Theta(\epsilon^{-1/3})\qquad(\epsilon\downarrow0)
\]
holds if and only if $x$ is badly approximable. If it is not, the stronger oscillation statements hold:
\begin{equation}\label{eq:r71-oscillation}
 \liminf_{\epsilon\downarrow0}\epsilon^{1/3}m_x(\epsilon)=0,
 \qquad
 \limsup_{\epsilon\downarrow0}\epsilon^{1/3}m_x(\epsilon)=+\infty.
\end{equation}
Thus even an existing logarithmic exponent $1/3$ need not give uniform cubic-root constants.
\end{theorem}
\begin{proof}
For badly approximable $x$, Theorem~\ref{thm:r69-cubic} already proves both uniform bounds, and that result is inherited. Otherwise Lemma~\ref{lem:r71-cf} supplies consecutive denominators $Q,Q'$ with $Q'/Q\to\infty$ along a subsequence. At $\epsilon'_n=e_x(Q)/Q$, the envelope is at most $Q$, and the upper in \eqref{eq:r71-envelope} gives
\[
 (\epsilon'_n)^{1/3}m_x(\epsilon'_n)
 \le C_0(Q^2e_x(Q))^{1/3}
 <C_0(Q/Q')^{1/3}\longrightarrow0.
\]
At the crossover noise \eqref{eq:r71-crossover}, the lower in \eqref{eq:r71-envelope} and \eqref{eq:r71-crossover-envelope} instead give
\[
 \epsilon_n^{1/3}m_x(\epsilon_n)
 \ge\tfrac14(Q'/Q)^{1/3}\longrightarrow+\infty.
\]
Both noise sequences tend to zero. Nonnegativity proves the first equality in \eqref{eq:r71-oscillation}, and the second rules out a uniform upper constant. No convergence of sampled continued-fraction coefficients has been substituted for a global arithmetic property.
\end{proof}

These conclusions distinguish three questions: the complete envelope at a given noise, the logarithmic exponents determined by $\mu$, and uniform power-law constants governed by bounded convergent ratios. A quadratic irrational phase has bounded partial quotients and retains the inherited cubic-root law. For an explicit phase with exponent two but without uniform cubic constants, take $x=[0;1,2,3,4,\ldots]$. Its denominator ratios are unbounded, while the denominator recursion implies at least Fibonacci growth and
$\log(q_{n+1}/q_n)/\log q_n\le\log(n+2)/\log q_n\to0$.
Hence $\mu=2$ but \eqref{eq:r71-oscillation} holds. For $\mu=\infty$, the two exponents are $0$ and $1/2$. The stronger universal theorem does not determine the numerical value of $\mu$ for the distinct algebraic unit phase $(3+4i)/5$; no cubic-root claim for that phase is introduced.

\subsection{Several visible frequencies and leakage}
Let $r\ge1$, let $x_j\in(0,1/2)$ be distinct, and let $w_j>0$ with $\sum_jw_j\le1$. For prescribed phases $\varphi_j$, consider
\[
 b_n^{\rm mix}=\tfrac12\left[1+(1-\epsilon)^n
              \sum_{j=1}^r w_j\cos(2\pi nx_j+\varphi_j)\right].
\]
Let $m_{\rm mix}$ be its least exact live memory. For $0<\lambda<1$, let $s_{\rm mix}$ be the least memory of the decaying response $\lambda^n b_n^{\rm mix}$.

\begin{corollary}[Multi-frequency arithmetic envelopes]\label{cor:r71-mixtures}
For every $0<\epsilon<1$,
\begin{align}
 \max\{2r+1,\tfrac12\max_j\mathfrak A_{x_j}(\epsilon)\}
 &\le m_{\rm mix}(\epsilon)
 \le2+C_0\sum_j\mathfrak A_{x_j}(\epsilon),\label{eq:r71-mixture}\\
 \max\{2r+2,1+\tfrac12\max_j\mathfrak A_{x_j}(\epsilon)\}
 &\le s_{\rm mix}(\epsilon)
 \le3+C_0\sum_j\mathfrak A_{x_j}(\epsilon).\label{eq:r71-leaky-mixture}
\end{align}
In particular each count is comparable, with constants depending only on $r$, to the largest of the component envelopes. If at least one phase is irrational, their upper logarithmic exponent is
\[
 \max_{j:\,x_j\notin\mathbb Q}
       \frac{\mu(x_j)-1}{2\mu(x_j)-1},
\]
with the same infinite-exponent convention. We do not replace the lower limit of a maximum by the maximum of lower limits; good approximation scales for different phases need not coincide. For a single leaking frequency, the exact identity $s=1+m_x$ transfers both exponents and the uniform-law rigidity unchanged, independently of $\lambda$.

These are actual boundary-optimal squared-task observer counts under the retained embedding. The protected target can be realized on dimension at most $2r+1$ before leakage and at most $2r+2$ after leakage; the flag embedding doubles these sufficient physical dimensions. No optimal physical-dimension claim or uniform controller for unknown target phases is made.
\end{corollary}
\begin{proof}
The distinct modes $1,(1-\epsilon)e^{\pm2\pi ix_j}$ have nonzero coefficients. Their rank and individual spectral obstructions, together with the counted block-mixture construction, are exactly those of Theorem~\ref{thm:r69-finite-mixture}. Apply \eqref{eq:r71-envelope} to each component to obtain \eqref{eq:r71-mixture}.

For leakage, the $2r+1$ decaying scalar modes gain the independent constant mode in the complementary output, giving rank $2r+2$. The relative-pole obstruction of Proposition~\ref{prop:r70-relative-spectrum} requires $1+m(1-\epsilon,2\pi x_j)$ labels for each $j$. Conversely append one live zero-output state to the entire unattenuated stochastic realization and scale its continuing matrix by $\lambda$. This realizes the decaying mixture from the same initial label. These arguments prove \eqref{eq:r71-leaky-mixture}.

For a finite list, the upper logarithmic limit of its maximum is the maximum of its upper logarithmic limits. Rational components have bounded exact size and do not change this maximum when an irrational component is present. Theorem~\ref{thm:r71-exponents} proves the formula. The single-frequency assertion follows from the exact leakage increment of Theorem~\ref{thm:r70-leakage}. Finally the block-diagonal protected realization in Theorem~\ref{thm:r69-finite-mixture}, with one additional leakage state when needed, supplies the stated dimensions; the actual squared-loss and stopping-cost assertions follow from the same declared flag embedding. A private randomized block choice is stored in the counted labels and never reveals the protected block.
\end{proof}

The conclusions do not supply a sharp leading constant, a general exact intermediate-memory risk curve, or a minimal controller for an unspecified noncomputable phase. For a supplied computable phase and positive noise, interval bounds on $e_x(q)$ and the explicit polygon construction give finite certified upper constructions; exact minimality still uses the separate spectral-membership problem. Ideal finite labels, descriptions of transition probabilities, exact sampling and working storage remain distinct resources. The complete arithmetic envelope closes the fixed-phase constant-factor question for this response family, not the general multi-command foundations programme.
